\vkthmsync{19}
\paragraph{\trDirected{} A∞-\trCategory。}\label{sub:directed-categories}
直到现在为止，我们采取的都是一种外在的观点：在一个取定的\trTriangulated{}
\trACategory{}的内部研究\trExceptionalCollection{}及其\trMutation。另一种做法
是，把注意力限制在一类特殊的 \trACategory{}上，这类\trCategory{}的\trObject{}%
自身就在\trDerivedCategory{}中构成一个\trExceptionalCollection。在这种内在的
观点之下，\trMutation{}将是一种把这样一个\trCategory{}变成另一个的操作，而后者
的\trDerivedCategory{}与前者\trEquivalence。

一个\textbf{\trDirected\trACategory}（directed A∞-category）是一个带有 $m$ 个%
\trObject{}的\trStrictlyUnital{} \trACategory{} $\scr{A}^{\rightarrow}$，连同
这些\trObject{}的一个排序 $(Y_1,\ldots,Y_m)$，使得
\begin{equation*}
\hom_{\scr{A}^{\rightarrow}}(Y_i,Y_j)=
\begin{cases}
0 & \text{当 } i>j \text{ 时},\\
\KK\cdot e_{Y_i} & \text{当 } i=j \text{ 时},\\
\text{在 }\KK\text{ 上有限维} & \text{当 } i<j \text{ 时。}
\end{cases}
\end{equation*}
由于\trStrictlyUnital{}性的条件，唯一非平庸的\trComposition{}映射是
\begin{equation}\label{eq:directed-composition-maps}
\begin{aligned}
\mu^d_{\scr{A}^{\rightarrow}}\colon
&\hom_{\scr{A}^{\rightarrow}}(Y_{i_{d-1}},Y_{i_d})\otimes\cdots\otimes
\hom_{\scr{A}^{\rightarrow}}(Y_{i_0},Y_{i_1})\\
&\longrightarrow\hom_{\scr{A}^{\rightarrow}}(Y_{i_0},Y_{i_d})[2-d],
\quad i_0<\cdots<i_d.
\end{aligned}
\end{equation}
因此，$\scr{A}^{\rightarrow}$ 可以由 $\KK$ 中作为
\eqref{eq:directed-composition-maps} 的矩阵\trCoefficients{}出现的那有限多个元素
完全描述。
令 $\scr{A}$ 是 $\scr{A}^{\rightarrow}$ 的\trTriangulated\trEnvelope。按定义，
诸 $Y_k$ 在\trCohomology\trFullyFaithful{}的\trFunctor{}%
$\scr{A}^{\rightarrow}\to\scr{A}$ 下的\trImage{}构成一个%
\trExceptionalCollection；由于 $\scr{A}$ 按构造由 $\scr{A}^{\rightarrow}$ 的%
\trObject{}所生成，这个序列是完整的。

\begin{remark}\label{thm:directed-complex-simplifications}
让我们更明确地取 $\scr{A}=\mathit{Tw}\scr{A}^{\rightarrow}$。\trDirected{}性
使得我们能够对\trTwistedComplex{}的定义作出若干细小的技术简化。首先，任何%
一个这样的\trComplex{}都可以写成
\begin{equation}\label{eq:directed-tw-complex-form}
\bigl(X=\bigoplus_{i=1}^{d}V^i\otimes Y_i,\ \delta_X=(\delta_X^{ji})\bigr)
\end{equation}
其中 $V^1,\ldots,V^d$ 是一些有限维\trGraded\trVectorSpace，而 $\delta_X^{ji}$
是一些\trDegree{}为 $1$ 的元素：
$\delta_X^{ji}\in\Hom_{\KK}(V^i,V^j)\otimes\hom_{\scr{A}^{\rightarrow}}(Y_i,Y_j)$。
对角项形如 $\delta_X^{ii}=\partial_{V^i}\otimes e_{Y_i}$，其中
$(V^i,\partial_{V^i})$ 是一个\trVectorSpace\trChainComplex。

下一个观察是：$\delta_X$ 的下三角性条件会自动得到满足，因而在这个特定情形
下可以从\trTwistedComplex{}的定义中省去。具体地说，可以用
$\bigoplus_{i\ge k}V^i\otimes Y_i$\footnote{译注：原书此处第二个因子印作
$Y^i$（带上标）；疑应为 $Y_i$（与\eqref{eq:directed-tw-complex-form} 中
$V^i\otimes Y_i$ 的下标记法平行，$Y^i$ 在本节无定义）；此处径改。——译者注}
来过滤\eqref{eq:directed-tw-complex-form}，然后再利用前导\trChainComplex{}%
$(V^k,\partial_{V^k})$ 的\trDegree{}\trFiltration{}来细分它。这就给出了 $X$
的一个\trFiltration，相对于它，$\delta_X$ 满足所想要的下三角性条件。

最后，任何\trTwistedComplex{} $X$ 在 $H^0(\scr{A})$ 中都\trIsomorphic{}于某个
对角项为零的 $\widetilde{X}$，即 $\delta_{\widetilde{X}}^{ii}=0$。具体地说，
Postnikov 分解\eqref{eq:postnikov-ladder} 表明，在相差\trIsomorphism{}的
程度下，可以从 $X_0=0$ 出发、通过一列\trMappingCone{}重构出 $X$：
\begin{equation}\label{eq:directed-postnikov-cone}
X_k\cong\mathrm{Cone}\bigl(S^{-1}X_{k-1}\longrightarrow Z_k\otimes Y_k\bigr).
\end{equation}
假设我们有某个与 $X_{k-1}$ \trIsomorphic{}的\trObject{} $\widetilde{X}_{k-1}$，
它是一个仅由\trDirectSummand{} $Y_1,\ldots,Y_{k-1}$ 构成、并且其%
\trDifferential{}的对角项为零的\trTwistedComplex。令 $\widetilde{X}_k$ 是对应于
\eqref{eq:directed-postnikov-cone} 的\trMorphism{}%
$S^{-1}\widetilde{X}_{k-1}\to Z_k\otimes Y_k$ 的\trCone。它与 $X_k$
\trIsomorphic，并且具有与 $\widetilde{X}_{k-1}$ 相同的性质，于是迭代这一构作，
最终便得到所要的 $\widetilde{X}=\widetilde{X}_m$。
\end{remark}
